\documentclass{article}

\usepackage{amsmath,amssymb}
\usepackage{xurl}
\usepackage[hidelinks]{hyperref}

% Shared commands for notes in docs/paper-gaps/.

\DeclareUnicodeCharacter{2080}{\ifmmode{}_{0}\else\textsubscript{0}\fi}
\DeclareUnicodeCharacter{2081}{\ifmmode{}_{1}\else\textsubscript{1}\fi}
\DeclareUnicodeCharacter{2082}{\ifmmode{}_{2}\else\textsubscript{2}\fi}
\DeclareUnicodeCharacter{2099}{\ifmmode{}_{n}\else\textsubscript{n}\fi}

\newcommand{\Lean}{\textsc{Lean}}
\ExplSyntaxOn
\NewDocumentCommand{\leanid}{m}{
  \ifmmode
    \textsf{\detokenize{#1}}
  \else
    \group_begin:
    \urlstyle{sf}
    \tl_set:Nn \l_tmpa_tl {#1}
    \tl_replace_all:Nnn \l_tmpa_tl {\_} {_}
    \exp_args:NV \path \l_tmpa_tl
    \group_end:
  \fi
}
\ExplSyntaxOff
\providecommand{\tr}{\operatorname{tr}}
\newcommand{\tnleanIssueBase}{https://github.com/LionSR/TNLean/issues/}
\newcommand{\tnleanPrBase}{https://github.com/LionSR/TNLean/pull/}
\newcommand{\tnleanDocBase}{https://sirui-lu.com/TNLean/docs/}
\newcommand{\ghissue}[1]{\href{\tnleanIssueBase#1}{GitHub issue \##1}}
\newcommand{\ghpr}[1]{\href{\tnleanPrBase#1}{PR \##1}}
\newcommand{\doclink}[1]{\href{\tnleanDocBase#1.html}{\path{#1}}}

% Dirac notation and the identity operator, as in blueprint/src/macros/common.tex.
% \providecommand keeps note-local definitions authoritative.
\usepackage{dsfont}
\providecommand{\ket}[1]{|#1\rangle}
\providecommand{\bra}[1]{\langle #1|}
\providecommand{\braket}[2]{\langle #1 | #2 \rangle}
\providecommand{\ketbra}[2]{|#1\rangle\!\langle #2|}
\providecommand{\Id}{\mathds{1}}
\providecommand{\one}{\mathds{1}}
\providecommand{\C}{\mathbb{C}}
\providecommand{\MN}[1]{M_{#1}(\C)}
\providecommand{\MD}{M_D(\C)}

% External referencing for paper sources.  The build helper writes sanitized
% auxiliary files under build/paper-aux/ and excludes duplicated source labels.
\usepackage{xr}

% Import a generated paper auxiliary file with a caller-supplied label prefix.
% The candidate paths support builds from docs/paper-gaps/, the repository root,
% and build/paper-aux/ itself.
\newcommand{\importpaperaux}[2]{%
  \IfFileExists{../../build/paper-aux/#2.aux}{%
    \externaldocument[#1]{../../build/paper-aux/#2}%
  }{%
    \IfFileExists{build/paper-aux/#2.aux}{%
      \externaldocument[#1]{build/paper-aux/#2}%
    }{%
      \IfFileExists{#2.aux}{%
        \externaldocument[#1]{#2}%
      }{}%
    }%
  }%
}

% General external reference macros
\newcommand{\paperref}[2]{\ref{#1-#2}}
\newcommand{\papereqref}[2]{(\ref{#1-#2})}

% CPSV16 source (1606.00608 / MPDO-22-12-17-2.tex) external document and shortcuts
\importpaperaux{cpsv16-}{1606.00608/MPDO-22-12-17-2-tnlean}
\newcommand{\cpsvref}[1]{\paperref{cpsv16}{#1}}
\newcommand{\cpsveqref}[1]{\papereqref{cpsv16}{#1}}

% Verdict marker: \gapnote{<kind>}{<status>}. Kind is the policy
% classification (clarification, local-correction, scope-restriction,
% unfaithful, false-source, open-gap); status is open, wip (elimination
% underway), resolved, or historical. Machine-read by the site index;
% prints a verdict line.
\newcommand{\gapnote}[2]{\par\noindent\textbf{Verdict:} #1 (#2).\par}


\title{The Sign of $\sigma_y$ in the Pauli Form of the AKLT Tensor}
\author{}
\date{2026-09-25}

\begin{document}
\maketitle
\gapnote{local-correction}{resolved}

\section{The source statement}

The examples appendix of Ref.~\cite{Cirac2021Matrix}, subsection ``The AKLT
state'', prints the AKLT tensor with the physical states labelled by
$S_z=0,\pm1$ as
\begin{align}
    A^{+1}
        &=\begin{pmatrix}1&0\\0&0\end{pmatrix}Y,
    &
    A^{0}
        &=\frac{1}{\sqrt2}\begin{pmatrix}0&1\\1&0\end{pmatrix}Y,
    &
    A^{-1}
        &=\begin{pmatrix}0&0\\0&1\end{pmatrix}Y,
    \label{eq:rmp_aklt_pauli_basis_sign_tensor}
\end{align}
with the singlet matrix
\begin{align}
    Y
        &=\begin{pmatrix}0&-1\\1&0\end{pmatrix}.
    \label{eq:rmp_aklt_pauli_basis_sign_singlet}
\end{align}
It then states that in the basis
\begin{align}
    \ket{+}
        &=\frac{i}{\sqrt2}\left(\ket{-1}+\ket{+1}\right),
    &
    \ket{-}
        &=\frac{1}{\sqrt2}\left(\ket{-1}-\ket{+1}\right),
    &
    \ket{0},
    \label{eq:rmp_aklt_pauli_basis_sign_basis}
\end{align}
the tensor becomes $A^-=\sigma_x/\sqrt2$, $A^+=\sigma_y/\sqrt2$, and
$A^0=\sigma_z/\sqrt2$.\footnote{Source traceability:
\path{Papers/2011.12127/TN-Review-main.tex}, lines 2372--2393.}

\section{The computation}

Write $\ket{\psi}=\sum_{s}\operatorname{tr}(A^{s_1}\cdots A^{s_N})
\ket{s_1\cdots s_N}$. The basis
(\ref{eq:rmp_aklt_pauli_basis_sign_basis}) is orthonormal, so the component of
$\ket{\psi}$ along $\ket{a_1\cdots a_N}$ is obtained from the tensor
\begin{align}
    B^{a}
        &=\sum_{s}\langle a\,|\,s\rangle\,A^{s}.
    \label{eq:rmp_aklt_pauli_basis_sign_new_tensor}
\end{align}
The matrices in (\ref{eq:rmp_aklt_pauli_basis_sign_tensor}) are
$A^{+1}=-\ketbra{0}{1}$, $A^{0}=\sigma_z/\sqrt2$, and
$A^{-1}=\ketbra{1}{0}$. Since $\langle +|\pm1\rangle=-i/\sqrt2$,
\begin{align}
    B^{+}
        &=-\frac{i}{\sqrt2}\left(A^{-1}+A^{+1}\right)
        =-\frac{i}{\sqrt2}\begin{pmatrix}0&-1\\1&0\end{pmatrix}
        =-\frac{1}{\sqrt2}\sigma_y,
    \label{eq:rmp_aklt_pauli_basis_sign_plus}
\end{align}
while $B^-=(A^{-1}-A^{+1})/\sqrt2=\sigma_x/\sqrt2$ and
$B^0=\sigma_z/\sqrt2$ agree with the source. The printed $\sigma_y$ is the
value of $\sum_s\langle s\,|\,a\rangle A^s$, which uses the complex conjugate
of the correct coefficient; it is not a gauge artefact. No virtual gauge
maps $(\sigma_x,\sigma_y,\sigma_z)$ to $(\sigma_x,-\sigma_y,\sigma_z)$,
because that map reverses the orientation of the Pauli vector, and the two
tensors give different states: for $N=3$ the coefficients of
$\ket{-,+,0}$ are $\operatorname{tr}(\sigma_x\sigma_y\sigma_z)/2^{3/2}=i/\sqrt2$
and its negative.

\section{The correction}

The printed Pauli form holds exactly in the orthonormal basis
\begin{align}
    \ket{+}
        &=-\frac{i}{\sqrt2}\left(\ket{-1}+\ket{+1}\right),
    \label{eq:rmp_aklt_pauli_basis_sign_corrected}
\end{align}
with $\ket{-}$ and $\ket{0}$ unchanged; equivalently, in the printed basis the
component along $\ket{+}$ is $-\sigma_y/\sqrt2$. The two readings differ by the
sign of one physical basis vector, so every symmetry, normality, and
parent-Hamiltonian property of the Pauli form is unaffected. Both identities,
the orthonormality of both bases, and the expansion of the state coefficients
in the new basis are formally verified in
\doclink{TNLean/MPS/Examples/AKLTReview}.\footnote{In particular,
\leanid{MPSTensor.physicalBasisChange_akltRMPBasis},
\leanid{MPSTensor.physicalBasisChange_akltRMPBasisFixed}, and
\leanid{MPSTensor.mpv_physicalBasisChange}.}

\section{Verdict}

The printed Pauli form of the AKLT tensor in Ref.~\cite{Cirac2021Matrix} holds
with $A^+=-\sigma_y/\sqrt2$ in the printed basis, or with the phase of
$\ket{+}$ conjugated. The tensor itself, its derivation from singlet bonds,
and the statements about $A^-$ and $A^0$ hold as printed.

\bibliographystyle{alpha}
\bibliography{references}

\end{document}
